\section{Fichamento 2: McElheran et al. (2024)}

\noindent\textbf{Autoria do fichamento:} Pedro Dias Soares

\subsection*{1. Referência bibliográfica}
\noindent\fullcite{mcelheran2024}

\subsection*{2. Tema do artigo}
A adoção inicial de inteligência artificial (IA) pelas empresas dos Estados Unidos: quem adotou,
que tecnologias adotou e onde as empresas adotantes ficam, com base no primeiro levantamento
oficial e representativo sobre o assunto.

\subsection*{3. Problema investigado}
Faltava evidência sistemática e representativa sobre o uso efetivo de IA nas empresas (Introdução,
p.~375). As medidas disponíveis estimavam o \emph{potencial} da IA de forma indireta, por descrições
de tarefas, anúncios de vagas e patentes (p.~377), e as pesquisas privadas que mediam uso tinham
amostras enviesadas para empresas grandes e de capital aberto (p.~382). Sem o \emph{denominador}
certo, que inclui as muitas empresas pequenas e jovens, não dá para saber quão difundida a IA está
nem onde seus efeitos devem aparecer primeiro.

\subsection*{4. Objetivo do trabalho}
Descrever a difusão inicial de cinco tecnologias associadas à IA na economia americana e
identificar as características de empresas, donos, estratégias e locais que se associam ao uso.
O objetivo é descritivo, nas palavras dos autores: ``Our main objective is to document'' (p.~377)
fatos estilizados, ``rather than explain'' (p.~377).

\subsection*{5. Hipótese ou pressuposto}
O artigo não testa uma hipótese única, mas declara expectativas explícitas na Seção~2
(pp.~378--380): donos com mais escolaridade e experiência tenderiam a liderar empresas mais
intensivas em tecnologia; donos com motivação de \emph{estilo de vida} adotariam menos; e capital
de risco e capital inicial alto sinalizariam potencial de crescimento ligado à IA (todas na
p.~379).

Há também um pressuposto que reconstruí: quem adota cedo antecipa quem vai adotar depois, de modo
que um retrato de 2017 indica a trajetória futura da IA. Ele aparece quando os autores falam dos
adotantes iniciais ``and, by extension, the most likely later'' (p.~377) e quando dizem que o
caminho da IA dependerá da sua dispersão geográfica inicial (p.~380). Não é testado, porque os
dados são de um único ano.

\subsection*{6. Base de dados ou corpus}
O levantamento é o \emph{Annual Business Survey} (ABS) de 2018, do Census Bureau com a NSF, com
referência ao ano de 2017. Foi enviado a 850.000 empresas, teve taxa de resposta de 69\% e cerca de
583.000 respondentes (pp.~376 e~380). As respostas foram ligadas à \emph{Longitudinal Business
Database} (LBD), base administrativa com emprego e receita ao longo da vida das empresas. A
amostra-base tem 447.000 empresas, que representam cerca de 4.050.000 com os pesos; a amostra com
dono e receita tem 209.000; e a de \emph{startups} de 0 a 5 anos, 75.000 (Tabela~1, p.~381). O mapa
usa cerca de 30.000 \emph{startups} das áreas metropolitanas com mais de um milhão de habitantes
(Fig.~5, p.~390).

Conta como uso de IA a resposta positiva para ao menos uma de cinco tecnologias: veículos guiados
automaticamente, aprendizado de máquina, visão computacional, processamento de linguagem natural e
reconhecimento de voz (Tabela~2, p.~382). Qualquer intensidade vale, inclusive \emph{menos de 5\% da
produção}.

\subsection*{7. Metodologia ou abordagem}
O estudo é observacional, descritivo e de corte transversal. A primeira parte (Seção~3,
pp.~380--390) traz estatísticas descritivas ponderadas: taxas de uso por tecnologia, porte, setor,
idade, características do dono e área metropolitana, com pesos de empresa, para representar a
população, e pesos de emprego, para medir a exposição dos trabalhadores (Figs.~1--5 e Tabelas~3
e~5).

A segunda parte (Seção~4, pp.~390--396) usa modelos de probabilidade linear por mínimos quadrados.
O uso de IA, uma variável 0/1, é regredido sobre porte, idade, digitalização e nuvem (Tabela~4) e,
para as \emph{startups}, sobre crescimento inicial da receita, características do dono,
financiamento, inovação e porte da área metropolitana (Tabela~7). Todos os modelos têm controles de
estado por indústria de seis dígitos, e cada coeficiente se lê como variação, em pontos percentuais,
da probabilidade de usar IA. Antes disso, os autores validam os marcadores de crescimento
regredindo o crescimento médio da receita nos três primeiros anos sobre eles (Tabelas~6a e~6b,
eq.~1).

\subsection*{8. Principais resultados}
Só 5,8\% das empresas usavam alguma das cinco tecnologias, e 1,1\% estavam testando; ponderado por
emprego, o uso sobe para 18,2\% (Fig.~1, p.~383). O aprendizado de máquina lidera, com 2,9\%, e os
veículos guiados ficam em 0,8\%. A maioria das empresas com mais de 5.000 empregados usa IA
(Fig.~2a). Indústria e informação ficam em torno de 12\%, e construção e varejo, em torno de 4\%
(Fig.~2b, p.~384). Usar nuvem se associa a 5 pontos percentuais a mais na probabilidade de usar IA
(0,050), e ter informação digitalizada, a 1,9 ponto (0,019) (Tabela~4, p.~386). Entre as empresas
que usam robôs, 57,5\% também usam IA (p.~385).

Entre as \emph{startups} (Tabela~7, coluna~5, pp.~395--396, $R^2 = 0{,}230$), os maiores
coeficientes são capital de risco (+8,6 pontos; 0,0862), proteção de propriedade intelectual
considerada \emph{muito importante} (+6,1), inovação de processo (+5,4) e patentes (+3,5). Inovação
de processo e propriedade intelectual, juntas, somam 11,5 pontos, quase o dobro da taxa média de 6\%
(p.~393). Um aumento de 10 pontos no crescimento inicial da receita eleva a probabilidade em 0,13
ponto na coluna~1 e em 0,06 na coluna~5.

Na geografia, os polos ficam acima da média no Vale do Silício e em metrópoles do Sul e do Oeste,
como Nashville, San Antonio, Las Vegas e Tampa (Fig.~5, pp.~389--390); com peso de emprego, entram
Louisville, Columbus, Austin e Atlanta (Fig.~A2). Estar numa área metropolitana grande eleva o uso
em 0,74 ponto (coluna~4, 0,0074), mas o efeito deixa de ser significativo quando entram as
características da empresa (coluna~5, 0,0042).

\subsection*{9. Conclusões dos autores}
O uso inicial de IA era baixo na média, mas concentrado em empresas grandes, inovadoras e com
marcadores de alto crescimento, e em alguns lugares (Seção~5, pp.~396--399). Dados e nuvem são
tecnologias habilitadoras centrais, e por isso os autores esperam um impacto econômico maior do que
a taxa média sugere. A conclusão diz que um possível \emph{AI divide} entre tipos de empresa e entre
lugares ``is visible'' (p.~399) e precisa ser avaliado, e reconhece que estabelecer causalidade
está ``beyond the reach of this first'' (p.~399) corte transversal de dados.

\subsection*{10. Limitações do estudo}
Os autores reconhecem (pp.~391 e~399) que a análise detalhada só vale para empresas jovens com dono
identificado, o que exclui as de capital aberto; que há viés de sobrevivência, porque empresas que
adotaram IA e fecharam antes de 2018 não aparecem; que não há causalidade; e que cerca de 80\% da
variação no uso de IA fica sem explicação (p.~399).

Percebi outras. A medida é frouxa e autodeclarada: conta como uso qualquer resposta positiva,
inclusive em \emph{menos de 5\%} da produção (Tabela~2 e nota da Tabela~7). A nota~2 (p.~400)
admite que os resultados refletem sobretudo \emph{machine learning}, e a definição desse termo na
Tabela~2 é ampla o bastante para abranger usos triviais de modelos preditivos. Os dados também são
antigos para o nosso tema: a referência é 2017, antes da IA generativa comercial, e a taxa de 5,8\%
(Fig.~1) não diz nada sobre o acesso a modelos como os de hoje.

O ganho de ajuste das características organizacionais é pequeno. A frase de que elas
``contributes significantly to the explanatory power'' (p.~399) se apoia num $R^2$ que passa de
0,209 para 0,230 (Tabela~7), ou 2,1 pontos absolutos, apresentados como aumento de ``about 10\%''
(p.~393). A maior parte do $R^2$ vem dos controles de estado por indústria, e o artigo não mostra um
modelo só com esses controles. As descritivas são ponderadas (p.~381), mas as notas das Tabelas~4
e~7 não dizem se as regressões usam os pesos; marco como \emph{[não informado]}. Por fim, a
geografia é só uma comparação entre metrópoles de um país: a Fig.~5 cobre áreas com mais de um
milhão de habitantes (nota~21, p.~401), sem intervalo de confiança por área. Para a pergunta do
nosso grupo, que compara países, o artigo serve de método, não de resultado.

\subsection*{11. Contribuições}
A contribuição principal é de medição. É a primeira estimativa representativa do uso de IA pelas
empresas de um país inteiro, com o denominador certo e ligada a dados administrativos de emprego e
receita (pp.~377 e~380--382). Mostra também que a taxa por empresa e a taxa por emprego contam
histórias diferentes (5,8\% contra 18,2\%, Fig.~1). Para o nosso projeto, ficam duas lições: a
escolha da unidade e do peso muda o tamanho da desigualdade medida, e a desigualdade geográfica
bruta pode ser composicional, já que o efeito do porte da área metropolitana some quando entram as
características das empresas (Tabela~7, colunas~4 e~5).

\subsection*{12. Avaliação crítica}
A amostra é muito grande e representativa, com alta taxa de resposta (69\%, p.~376), e vem de uma
fonte oficial ligada a registros administrativos, o que evita o viés das pesquisas privadas
(p.~382). Os autores são cuidadosos com causalidade e repetem que os achados são correlações
(pp.~391, 396 e~399 e nota~28). Também separam o peso por empresa do peso por emprego, o que é raro.

\textit{(a) Auditoria da tabela principal.} A afirmação central sobre quem adota se apoia na
Tabela~7 (pp.~395--396). A tabela traz erros-padrão, mas a soma de coeficientes que o texto destaca
(11,5 pontos para inovação de processo mais propriedade intelectual, p.~393) aparece sem erro-padrão
próprio. Por ser tabela, a questão do eixo vertical não se aplica. As colunas~1--5 têm 75.000
observações, número arredondado por sigilo, mas os coeficientes de capital de risco dependem de uma
fração pequena delas, porque só 0,9\% das \emph{startups} têm esse financiamento (Tabela~5, p.~388).
Falta um modelo só com os controles de estado por indústria, que mostraria quanto do $R^2$ de 0,209
a 0,230 vem das variáveis de interesse. Também falta a taxa de uso prevista para cada perfil de
empresa, que o texto só apresenta em termos relativos, como ``a 17\% increase at the mean''
(p.~393). Eu esperaria ver a coluna de base só com controles e as somas de coeficientes destacadas
no texto com seus intervalos de confiança.

\textit{(b) Exame de uma afirmação do resumo.} O resumo afirma que a adoção ``was far from evenly
distributed'' (p.~375) e que um punhado de cidades e polos emergentes liderou a adoção pelas
\emph{startups}. A peça principal é a Fig.~5 (p.~390), e ela mostra posições, não taxas: a cor de
cada área é o percentil da sua taxa de uso entre as grandes metrópoles. Uma escala de percentis vai
sempre de 0 a 100, mesmo que as taxas sejam quase iguais, e a figura não mostra a taxa em si nem a
incerteza de cada área. As outras peças apontam diferenças pequenas. O próprio texto diz que as
áreas menores ficam ``not too far below'' (p.~390) as maiores; na Tabela~5, 55,1\% das
\emph{startups} que usam IA estão em grandes metrópoles, contra 49,9\% das que não usam; e na
coluna~5 da Tabela~7 o efeito do porte da área deixa de ser significativo. Classificação: conclusão
mais ampla que o experimento permite. O artigo sustenta uma ordenação das metrópoles, não que a
desigualdade entre elas seja grande nem que seja distinguível de ruído.
